\subsection{平衡点、偏差変数、線形化}

ロボットアームをある姿勢の近くで動かす、熱板をある温度の近くに保つ、モータをある回転速度の近くで使う、という設計では、対象が取りうるすべての状態を同じ精度で表す必要はない。選んだ動作点の近くで、位置や温度が少しずれたときに変化率がどれだけ変わるか、入力を少し増やしたときに速度や温度の変化率がどれだけ変わるか、という局所的な入力効果の傾きがまず必要になる。制御器はこの傾きを使って、小さなずれを小さな入力変更で打ち消す。

最初は状態も入力も 1 つだけの系
\begin{equation}
  \dot{x}=f(x,u)
\end{equation}
で考える。基準にしたい値を $x_e$、その値を保つ一定入力を $u_e$ とする。もしその点で状態が止まるなら
\begin{equation}
  f(x_e,u_e)=0
  \label{eq:scalar-equilibrium-condition}
\end{equation}
である。実際の状態と入力がこの基準から少しずれているとき、
\begin{equation}
  x=x_e+\tilde{x},
  \qquad
  u=u_e+\tilde{u}
\end{equation}
と書く。$x_e$ は一定値なので $\dot{x}=\dot{\tilde{x}}$ であり、元の式へ代入すると
\begin{equation}
  \dot{\tilde{x}}=f(x_e+\tilde{x},u_e+\tilde{u})
\end{equation}
となる。ここで基準点では \eqref{eq:scalar-equilibrium-condition} が成り立つため、同じ 0 を引いて
\begin{equation}
  \dot{\tilde{x}}
  =f(x_e+\tilde{x},u_e+\tilde{u})-f(x_e,u_e)
\end{equation}
と書ける。この差は、点 $(x_e,u_e)$ の近くでグラフを接線で見ることに対応する。そこで局所的な傾きを
\begin{equation}
  a=\left.\frac{\pp f}{\pp x}\right|_{(x_e,u_e)},
  \qquad
  b=\left.\frac{\pp f}{\pp u}\right|_{(x_e,u_e)}
\end{equation}
と置けば、小さい $\tilde{x}$、$\tilde{u}$ に対して
\begin{equation}
  f(x_e+\tilde{x},u_e+\tilde{u})-f(x_e,u_e)
  \simeq a\tilde{x}+b\tilde{u}
\end{equation}
である。したがって、動作点近傍の偏差は
\begin{equation}
  \dot{\tilde{x}}\simeq a\tilde{x}+b\tilde{u}
  \label{eq:scalar-linearization}
\end{equation}
で近似できる。ここで $a$ は状態のずれが変化率に現れる強さ、$b$ は入力変更が変化率へ効く強さである。非線形性を消したのではなく、選んだ点の近くで一次の傾きだけを残している。

この基準点とずれの考えを、多変数の状態方程式で使えるように定義として固定する。

\begin{definition}[平衡点]
時間に陽に依存しない系 $\dot{x}=f(x,u)$ において、一定入力 $u_e$ と状態 $x_e$ が
\begin{equation}
  f(x_e,u_e)=0
\end{equation}
を満たすとき、$(x_e,u_e)$ を平衡点という。
\end{definition}

平衡点では状態の変化率が 0 なので、入力を $u_e$ に保てば状態は $x_e$ に留まる。熱系なら一定温度、RC 回路なら一定電圧、モータなら一定速度、水槽なら一定水位が平衡点の例である。外乱やパラメータを明示するモデル
\begin{equation}
  \dot{x}=f(x,u,d,p)
\end{equation}
では、一定外乱 $d_e$ とパラメータ $p$ を固定したうえで
\begin{equation}
  f(x_e,u_e,d_e,p)=0
\end{equation}
を満たす組を平衡点として扱う。周囲温度や負荷トルクが変われば、同じ入力でも平衡点は変わりうる。時間に陽に依存するモデルや周期的に動く基準運動では、厳密には一点の平衡点ではなく基準軌道のまわりで偏差を考えるが、本節ではまず一定値に留まる動作点を扱う。

\begin{definition}[偏差変数]
制御では、絶対値そのものと同時に平衡点からのずれを見る。そこで
\begin{equation}
  \tilde{x}=x-x_e,
  \qquad
  \tilde{u}=u-u_e
  \label{eq:deviation-definition}
\end{equation}
を偏差変数という。すると $x=x_e+\tilde{x}$、$u=u_e+\tilde{u}$ である。
\end{definition}

偏差変数は無次元の誤差ではない。$x_i$ が温度なら $\tilde{x}_i$ も温度差の単位 K をもち、$u_j$ が力なら $\tilde{u}_j$ も N をもつ。必要なら後で基準値や許容幅で割って無次元化するが、線形化そのものでは物理単位を保ったまま計算する。平衡点は定数なので $\dot{x}_e=0$ であり、
\begin{equation}
  \dot{\tilde{x}}=\frac{\dd}{\dd t}(x-x_e)=\dot{x}
\end{equation}
である。出力も同じ考えで、$y_e=g(x_e,u_e)$ と置けば $\tilde{y}=y-y_e$ を出力偏差として使える。

状態が複数あると、式 \eqref{eq:scalar-linearization} の傾き $a$、$b$ は 1 つの数では足りない。各状態の変化率 $f_i$ が、各状態 $x_j$ と各入力 $u_j$ にどれだけ反応するかをすべて並べる必要がある。この傾きの表がヤコビアンである。

\begin{definition}[ヤコビアン]
ベクトル値関数 $f(x,u)\in\R^n$ の各成分を各変数で偏微分して並べた行列をヤコビアンという。$x\in\R^n$、$u\in\R^m$ なら
\begin{equation}
  \frac{\pp f}{\pp x}=\left[\frac{\pp f_i}{\pp x_j}\right]_{i,j},
  \qquad
  \frac{\pp f}{\pp u}=\left[\frac{\pp f_i}{\pp u_j}\right]_{i,j}
\end{equation}
である。
\end{definition}

$\dot{x}=f(x,u)$ では $f_i$ の単位は $x_i$/s である。したがって $\pp f_i/\pp x_j$ の単位は $(x_i/\mathrm{s})/x_j$、$\pp f_i/\pp u_j$ の単位は $(x_i/\mathrm{s})/u_j$ になる。$A=\pp f/\pp x$ は $n\times n$ 行列、$B=\pp f/\pp u$ は $n\times m$ 行列であり、$A\tilde{x}$ と $B\tilde{u}$ の各成分はどちらも $x_i$/s の単位をもつ。状態の各成分が位置、速度、電流のように違う単位をもつとき、行列の成分も同じ単位ではない。

\begin{theorem}[一次テイラー展開による線形化]
非線形系 $\dot{x}=f(x,u)$ が平衡点 $(x_e,u_e)$ の近くで十分なめらかなら、偏差変数 $\tilde{x}$、$\tilde{u}$ は一次近似で
\begin{equation}
  \dot{\tilde{x}}=A\tilde{x}+B\tilde{u}
  \label{eq:linearized-model}
\end{equation}
に従う。ただし
\begin{equation}
  A=\left.\frac{\pp f}{\pp x}\right|_{(x_e,u_e)},
  \qquad
  B=\left.\frac{\pp f}{\pp u}\right|_{(x_e,u_e)}
\end{equation}
である。
\end{theorem}

ここで「十分なめらか」とは、少なくとも平衡点の近傍で $f$ の一次偏微分が存在して連続であり、二次の誤差項を小さくできるという意味である。より安全には $f$ が二回連続微分可能で、二階偏微分が近傍で有界であると仮定する。この仮定により、$\delta z=[\tilde{x}^\T,\tilde{u}^\T]^\T$ と置いたとき、捨てた項 $r_2(\delta z)$ は
\begin{equation}
  \lVert r_2(\delta z)\rVert \le c\lVert \delta z\rVert^2
\end{equation}
を満たす定数 $c$ で抑えられる。したがって一次モデルは $\tilde{x}$ と $\tilde{u}$ が小さい範囲でだけ信頼できる。入力飽和、摩擦の符号切替、接触の有無のように式が滑らかでない現象をまたぐと、この線形化の仮定は崩れる。

\begin{derivation}
$x=x_e+\tilde{x}$、$u=u_e+\tilde{u}$ を $\dot{x}=f(x,u)$ に代入する。$x_e$ は定数なので $\dot{x}=\dot{\tilde{x}}$ である。したがって
\begin{equation}
  \dot{\tilde{x}}=f(x_e+\tilde{x},u_e+\tilde{u})
\end{equation}
となる。平衡点の式 $0=f(x_e,u_e)$ を右辺から引いても値は変わらないので
\begin{equation}
  \dot{\tilde{x}}
  =f(x_e+\tilde{x},u_e+\tilde{u})-f(x_e,u_e)
\end{equation}
である。多変数のテイラー展開を一次まで使うと
\begin{align}
  f(x_e+\tilde{x},u_e+\tilde{u})-f(x_e,u_e)
  &=
    \left.\frac{\pp f}{\pp x}\right|_{(x_e,u_e)}\tilde{x}
    +\left.\frac{\pp f}{\pp u}\right|_{(x_e,u_e)}\tilde{u}
    +\text{二次以上の項}.
\end{align}
この消去は、選んだ $(x_e,u_e)$ が平衡点であることに依存する。平衡点でない点を基準にすると、テイラー展開には $f(x_e,u_e)$ という定数項が残り、偏差原点 $\tilde{x}=0$ は止まる点にならない。平衡点近傍だけを見る一次近似では二次以上の項を捨てるので
\begin{equation}
  \dot{\tilde{x}}\simeq A\tilde{x}+B\tilde{u}
\end{equation}
を得る。
\end{derivation}

\begin{example}[熱系の動作点まわりの線形化]
小さな物体の温度 $T(t)$ [K] を、周囲温度 $T_a$ [K] の環境でヒータ電圧 $v(t)$ [V] により調整する。熱容量を $C_{\mathrm{th}}$ [J/K]、熱抵抗を $R_{\mathrm{th}}$ [K/W]、ヒータの電圧二乗係数を $\eta$ [W/V$^2$] とし、ヒータ電力が $\eta v^2$ で近似できるとする。このとき
\begin{equation}
  C_{\mathrm{th}}\dot{T}
  =-\frac{T-T_a}{R_{\mathrm{th}}}+\eta v^2
\end{equation}
である。状態を $x=T$、入力を $u=v$ とすれば
\begin{equation}
  f(T,v)=\frac{1}{C_{\mathrm{th}}}
  \left(-\frac{T-T_a}{R_{\mathrm{th}}}+\eta v^2\right)
\end{equation}
である。一定電圧 $v_e$ で温度が止まる平衡点は
\begin{align}
  0&=-\frac{T_e-T_a}{R_{\mathrm{th}}}+\eta v_e^2,\\
  T_e&=T_a+R_{\mathrm{th}}\eta v_e^2
\end{align}
を満たす。偏差を $\tilde{T}=T-T_e$、$\tilde{v}=v-v_e$ と置くと、
\begin{align}
  \dot{\tilde{T}}
  &=\frac{1}{C_{\mathrm{th}}}
    \left(-\frac{T_e+\tilde{T}-T_a}{R_{\mathrm{th}}}
    +\eta(v_e+\tilde{v})^2\right)\\
  &=\frac{1}{C_{\mathrm{th}}}
    \left(\underbrace{-\frac{T_e-T_a}{R_{\mathrm{th}}}
    +\eta v_e^2}_{\text{平衡条件により }0}
    -\frac{\tilde{T}}{R_{\mathrm{th}}}
    +2\eta v_e\tilde{v}
    +\eta\tilde{v}^2\right).
\end{align}
平衡点を選んだ効果は、上の括弧内の定数項が消えることに現れる。残るのは偏差に比例する項と、入力偏差の二乗項であり、
\begin{equation}
  \dot{\tilde{T}}
  =-\frac{1}{R_{\mathrm{th}}C_{\mathrm{th}}}\tilde{T}
   +\frac{2\eta v_e}{C_{\mathrm{th}}}\tilde{v}
   +\frac{\eta}{C_{\mathrm{th}}}\tilde{v}^2
\end{equation}
となる。一次線形化では最後の二次項を捨てて
\begin{equation}
  \dot{\tilde{T}}
  \simeq
  -\frac{1}{R_{\mathrm{th}}C_{\mathrm{th}}}\tilde{T}
  +\frac{2\eta v_e}{C_{\mathrm{th}}}\tilde{v}
\end{equation}
を得る。ここで
\begin{equation}
  A=-\frac{1}{R_{\mathrm{th}}C_{\mathrm{th}}},
  \qquad
  B=\frac{2\eta v_e}{C_{\mathrm{th}}}
\end{equation}
であり、$A$ の単位は 1/s、$B$ の単位は K/(s\,V) である。$v_e=0$ の動作点では $B=0$ となり、電圧の小さな正負の偏差は一次近似では温度を変えないという式になる。これはヒータ電力を電圧の二乗で近似したためであり、制御入力を電力 $P=\eta v^2$ として選ぶか、実際に使う非零の動作点まわりで線形化する必要がある。
\end{example}

倒立振子で $x=[\theta,\dot{\theta}]^\T$ とし、
\begin{equation}
  f(x,u)=\begin{bmatrix}x_2\\ (mgl\sin x_1+bu)/J\end{bmatrix}
\end{equation}
と書く。直立平衡点 $(x_e,u_e)=(0,0)$ まわりのヤコビアンは
\begin{align}
  A&=\left.\frac{\pp f}{\pp x}\right|_{(0,0)}
   =\begin{bmatrix}
      \pp f_1/\pp x_1 & \pp f_1/\pp x_2\\
      \pp f_2/\pp x_1 & \pp f_2/\pp x_2
    \end{bmatrix}_{(0,0)}\\
   &=\begin{bmatrix}
      0&1\\
      (mgl/J)\cos x_1&0
    \end{bmatrix}_{x_1=0}
    =\begin{bmatrix}0&1\\ mgl/J&0\end{bmatrix},\\
  B&=\left.\frac{\pp f}{\pp u}\right|_{(0,0)}
   =\begin{bmatrix}0\\ b/J\end{bmatrix}.
\end{align}
この計算で、$\sin\theta\simeq\theta$ という近似を単独の経験則として使ったのではない。ヤコビアンを計算した結果として、直立近傍の一次モデルが得られている。
